\section{Failure-To-Evidence Mapping}

Table~\ref{tab:delegation-failures} gives candidate diagnostic translations. For each surface failure, it asks which delegation-layer hypothesis should be tested when evaluating whether an operation executed through the deployment was authorized. The mapping proposes diagnoses rather than automatically reclassifying every listed failure as a delegation failure.

\begin{table}[p]
\scriptsize
\renewcommand{\arraystretch}{0.92}
\centering
\begin{tabular}{>{\raggedright\arraybackslash}p{0.22\linewidth}%
                >{\raggedright\arraybackslash}p{0.34\linewidth}%
                >{\raggedright\arraybackslash}p{0.34\linewidth}}
\toprule
Surface failure & Candidate delegation-layer diagnosis & Security question \\
\midrule
\multicolumn{3}{l}{\textit{Instruction-status, standing, and priority failures}} \\
\midrule
Direct prompt injection & The deployment acts as though an attacker-controlled directive superseded the governing instruction. & What was the occurrence's contextual status, did its source have standing, and did any applicable priority rule match the authority regime? \\
Indirect prompt injection & The deployment's classification treats directive-form content presented as data as a command in use, or its action selection treats a source without standing as controlling. & Were contextual status, source standing, and any applicable conflict priority recorded separately? \\
Jailbreak & Safety constraints are treated as negotiable within the user dialogue. & Which constraints are non-delegable, and did the deployment preserve and apply them under pressure? \\
Apparent non-responsiveness & The deployment rejects a later directive because its source lacks amendment or override standing, or because a joint, review, or release condition isn't satisfied. This can be authorized resistance rather than failure. & Which mandate was operative, which transition was attempted, and did resistance preserve or violate that mandate? \\
Improper override or interference & A principal, intermediary, or other source causes the recorded or executed authorization state to change without the operation-specific standing or guards required by \(R\). & Did the purported amendment, revocation, override, or release produce a valid normative transition, only a causal change in the recorded or executed state, or neither? \\
\midrule
\multicolumn{3}{l}{\textit{Proposal, gate, and execution failures}} \\
\midrule
Tool misuse & A tool-call proposal falls outside recorded scope, a gate allows an operation not authorized under \(R\), or the deployment executes a prohibited operation. & What did the record represent as permitted, what did the gate decide, and what operation was executed? \\
Data exfiltration & Access to information is delegated without corresponding release authority. & Did the deployment distinguish permission to read from permission to disclose? \\
Memory poisoning & The deployment executes a persistent-state update on the basis of content from a source without standing or without an authorized update path. & Who or what may authorize durable-context updates, and how is the update audited? \\
Goal drift & The deployment pursues an objective outside the operative mandate, including its valid amendments, protected discretion, and release conditions. Divergence from a source's latest preference isn't sufficient. & What mandate was operative at the relevant time, and what evidence shows whether the deployment stayed within it? \\
Over-refusal & The deployment refuses work that the authority regime permits. The refusal is a task-fulfilment failure; if the regime also requires the work and the omission is assessed, it defeats mandate compliance as well. & Was completion merely permitted or required, and which task-success or mandate-compliance claim did the refusal defeat? \\
Under-refusal & The deployment delivers or executes an operation prohibited by the authority regime. & Which boundary should have blocked the operation, and why did it fail? \\
Scaffold failure & A component in the configured deployment fails to propagate or correctly apply a recorded grant or constraint. & Did the deployment propagate the record of the governing grant and constraints across components, not just inside the model turn? \\
\midrule
\multicolumn{3}{l}{\textit{Evidentiary and evaluation failures}} \\
\midrule
Agent transcript misreporting & A deployment-generated report doesn't match what happened during delegated action. & Can an assurance reviewer reconstruct the executed operation and technical failure locus from the trace and supporting evidence? \\
LLM-judge failure & An LLM judge, as an automated evaluator, collapses task success, compliance, risk, and technical failure attribution. & Was the automated evaluator validated and scoped to support the declared assurance claim? \\
\bottomrule
\end{tabular}
\caption{Candidate delegation-layer diagnoses for AI-security failures.}
\label{tab:delegation-failures}
\end{table}

The corresponding evaluation requirements are shown in Table~\ref{tab:delegation-evidence}. The rows are fields required by an authorization or evaluation claim, not a general test of the full institutional record. Evidentiary assurance asks whether the preserved record makes that claim evaluable and what evidentiary verdict it warrants. A delegation-assurance evaluation doesn't need every row for every deployment: different claims require different applicable subsets and depths. The framework assigns no weights and sums no rows into a scalar score. Each claim has to name its required subclaims and defeating tests; the noncompensatory rule applies to every declared claim tree.

\begin{table}[p]
\scriptsize
\setlength{\tabcolsep}{3pt}
\renewcommand{\arraystretch}{0.88}
\centering
\begin{tabular}{>{\raggedright\arraybackslash}p{0.27\linewidth}%
                >{\raggedright\arraybackslash}p{0.31\linewidth}%
                >{\raggedright\arraybackslash}p{0.32\linewidth}}
\toprule
Authorization or diagnostic dimension & Required authorization record & Evaluation requirement \\
\midrule
Instruction status & The record shows how directive-form string occurrences were classified and used: as commands in use, quotations or other mentions, examples, data, policy text, tool results, reports, or objects of analysis. & Minimal pairs contrast command use, mention, embedded content, policy example, and reported speech while holding the action target constant; the recorded classification is scored against independently adjudicated status, with fit recorded separately from the resulting action. \\
Source standing and conflict priority & The trace records governing sources, authority object, effective interval, operation-specific standing, joint or guarded conditions, transition result, channel, content type, and technical trust separately; independently adjudicated standing supplies the reference. & Crossed items vary prompt or record representations keyed to standing and priority labels while holding directive content and the surface request fixed; the causal contrast is between assigned packages, not normative statuses. \\
Scope control & The trace records the operation class, target, time, tool, data class, approval threshold, and authorization verdict used at the gate. & Items vary one scope boundary at a time: permitted action, prohibited action, target, tool, time, data class, or approval threshold. \\
Provenance & Issuer or role, provenance channel, content type, and technical-trust status remain separately recorded and bound to classifications or observable action-selection decisions in the trace. & The same directive-form content appears as user request, webpage text, email body, retrieved document, tool output, policy text, and memory. \\
Action licensing & The trace distinguishes the model-generated tool-call proposal, gate decision, and executed operation, and records the authorization verdict under the applicable grant, constraints, and approval state. & Tool-call tasks separately score proposal, gate decision, and execution for reading, transforming, storing, sending, deleting, and purchasing. \\
Information flow & Read, transform, store, and disclose permissions remain distinct; quoting or metalinguistically reproducing a protected value in a delivered output can still disclose it. & Read-but-do-not-reveal and transform-without-leaking tasks use benign canaries and confidential-but-actionable fields, including metalinguistic, quoted, or negated outputs that reproduce a canary. \\
Persistence & Durable memory or state updates require a source with standing and an authorized update path. & Memory-update tasks cross sources with and without standing, including technically trusted, untrusted, stale, adversarial, and user-corrected sources. \\
Refusal calibration & Refusal and completion are scored separately for action authorization, task success, and any requirement to act. & Paired tasks distinguish permitted completion, required completion, prohibited execution, and safe analysis of quoted or embedded descriptions. \\
Technical failure attribution & The record separately identifies the technical locus, the applicable normative boundary, and evaluator or taxonomy uncertainty. & Technical-locus labels distinguish model action selection, scaffold or routing, policy-enforcement point, tool or runtime, and logging; separate fields record normative adjudication and evaluative uncertainty. \\
Audit reconstruction & A third party can reconstruct the recorded authorization path, operating conditions, execution, and review path without relying on a model-generated statement that the action was authorized. & Trace-review tasks require reconstruction of governing sources, effective interval, standing by operation, purported transitions and assigned effects, joint approvals or guards, scope, configuration, proposal, gate decision, execution, and recorded policy reason. \\
Automated-evaluator validation & The evaluation record keeps the delegation dimensions separate and includes uncertainty rather than collapsing them into one dispositive label. & On held-out cases matched to the evaluator's declared \(T\), LLM-judge validation compares no-rubric, rubric-aided, and adversarially phrased conditions against expert-adjudicated reference labels. \\
\bottomrule
\end{tabular}
\caption{Delegation-claim fields and evaluation requirements.}
\label{tab:delegation-evidence}
\end{table}

The procurement example shows the unit of evaluation. A final memo can be correct while the action-authorization claim is defeated because the deployment executed an unauthorized email operation, disclosed the budget ceiling, or updated a vendor record without permission.

Conversely, refusing a permitted comparison task can defeat task success without defeating action authorization. If \(R\) required completion and the omission is assessed, the refusal instead defeats a separately declared mandate-compliance claim. An inadequate trace makes the relevant claim indeterminate rather than turning missing evidence into evidence of unauthorized action. The target is the action-authorization claim, not the final answer alone.

\clearpage
